\chapter{Roots and Weights}

\section{Motivating example: \(\mathfrak{su}(2)\)}

\subsection{Projective representations and simply-connected groups}

representations as linear maps on vectors

when we say the representation of a group is in GL or the one of an algebra is in End, we are saying that the group or algebra is represented as linear maps on vectors.

when we measure we normalize with resoect to the inner product, we need the norm of the vectors and thus the inner product...

all of the vectors on the same ray (same direction) are to be considered in the same state. it is a one dimensional subspace of the vector space. if this is what carachterizes the system, we need to really care only about the projected representation on the projective space, which is the space of rays. in qm symmetries should be realized as projective representattions of the group, rather than as linear representations. but we can always lift a projective representation to a linear one, by going to the universal cover of the group. for example, the projective representations of SO(3) are the linear representations of SU(2). this is why we can use SU(2) to describe spin, even if spin is not a representation of SO(3).

the mathematical key is the one saying that the a finite dimensional projective representation of a lie group on a finite dimensional complex vector space lifts to a unitary linear representation the covering group \(\overline{G}\), which is the unique simply connected Lie group with the same lie algebra as \(G\). hard to prove, absolutely non trivial.

projective representations \(\equiv\) unitary linear representations of the covering group \(\equiv\) anti hermitian linear representations of the lie algebra \(\equiv\) linear representations of the complexified lie algebra

In fact, for the special case of \(G=SO(3)\) and \(\mathfrak{g} = \mathfrak{su}(2)\), this math-result, combined with the “axiomatic” definition of quantum mechanics using the (projective) Hilbert space, can be viewed as predicting half-integer spin particles (i.e., linear representations of \(\mathfrak{su}(2)\) which are not \(SO(3)\) reps) in the quantum world: they arise from projective representations of the rotation symmetry \(SO(3)\), which we do not have in classical mechanics.

    [details on axiomatics of qm]

\subsection{Ladder Operators and Spin Representation}

\paragraph{Complexification and Hermitian basis.} from properties of complexification we can take complex linear combinations of the generators of the algebra, and we can choose them in a way that they have nice properties. since the original generators are anti-hermitian, we can take linear combinations that are hermitian. so we can diagonalize them. thinking of the adjoint rep in that basis, which is an endomorphism of the algebra, we can see that the adj rep of the generators is also an endomorphism of the algebra, and thus we can diagonalize it as well, since it will remain hermitian. this property will extend to all the other complex finite dimensional representations, and thus they will have a basis in which the generators are hermitian and thus diagonalizable.

\paragraph{Eigenvalues of \(d(J_3)\) from ladder operators.} we can pick an on basis thanks to the usual unitarity trick which allows us to make the representation unitary, and thus the generators hermitian. in the following we prove that for irreducible representations the dimension of the subspace of eigenvectors of \(d(J_3)\) with a given eigenvalue is one, and that the eigenvalues are equally spaced.
the ladder operators are not hermitian. after definingthem through the commutation relations, we can study how the eigenvalues of \(d(J_3)\) change when we apply the ladder operators (it moves up and down in the different eigenvector subspaces, the state changes but it remains an eigenvector). we can see that they change by one unit, and thus the eigenvalues are equally spaced (and thus the entire lie algebra moves among the invariant eigenvectors subspaces). this dont work in irreducible representations, since there do not exist trivial invariant subspaces.
if the representation is finite, there must exist a maximum range of eigenvalues, and thus there must be a highest weight vector, which is an eigenvector of \(d(J_3)\) that is annihilated by the raising operator. same for the minimum value (annihilated by the lowering operator).thus the representation space is made of a finite direct sum of eigenspaces of \(d(J_3)\) with equally spaced eigenvalues, and the ladder operators move us among them.

\paragraph{The Casimir operators and range of \(J_3\)-eigenvalues.} to further restrict the ...
the casimir operator commutes with all the generators, thus it is a multiple of the identity in an irreducible representation. we can apply schurs lemma to infer that the eigenvalue of the casimir operator is a constant in an irreducible representation.
if we apply the obtained formulae to the basis of eigenvectors of \(d(J_3)\), we can find a relation between the eigenvalue of the casimir operator and the range of the eigenvalues of \(d(J_3)\). this allows us to find the possible values of the eigenvalues of \(d(J_3)\) and thus to classify all the irreducible representations. the proprierties we derive constrain the spacing of the eigenvalues to be integer or half-integer (and \(>0\)), and thus we can label the irreducible representations by a non-negative integer or half-integer \(j\), which is the maximum eigenvalue of \(d(J_3)\) (the highest weight). the dimension of the representation is then \(2j+1\).

\paragraph{Dimensionality of eigenspaces.} we want to argue that all the eigenspaces of \(d(J_3)\) are one-dimensional. the span of the maximum eigenvectors and consecutive applications of the lowering operator is an invariant subspace, thus it must be the entire representation space, since the representation is irreducible. thus all the eigenvectors of \(d(J_3)\) are obtained by applying the lowering operator to the maximum eigenvector, and thus they are all one-dimensional (we never get to the same eigenvalue twice, since the eigenvalues are equally spaced and we have a finite number of them).

\paragraph{Normaliation and numerical coefficients.} [...] we find how to completely specify the representation by giving the value of \(j\) and the normalization of the ladder operators. we can choose a normalization such that the action of the ladder operators on the eigenvectors is given by a simple formula, which is very convenient for computations.

\section{Reformulation with Roots and Weights}

 [...]

\section{Killing Form and Cartan Subalgebra}

For a general Lie algebra \(\mathfrak{g}\), we first need the analogue of the subalgebra spanned by \(J_3\) in \(\mathfrak{sl}(2,\mathbb{C})\)  which we choose to diagonalize.

\begin{definition}[Killing form]
    [...]
\end{definition}

It takes two elements of the algebra and gives a number, thus it is a bilinear form on the algebra. since we are taking a trace of endomorphisms, we now there exist a basis invariant way to define it, and thus it is invariant under the adjoint action of the algebra on itself. this is a very important property, since it allows us to use the killing form to define an inner product on the algebra, which is invariant under the adjoint action. this will allow us to define a notion of orthogonality and thus to define a notion of diagonalization for the adjoint representation, which is crucial for the classification of representations.

The symmetric property comes from the fact that the trace of a product of endomorphisms is invariant under cyclic permutations, thus we can swap the two elements without changing the value of the killing form. the invariance property comes from the fact that the killing form is defined in terms of the adjoint representation, which is invariant under the adjoint action.

Given a basis for the lie algebra, given by the generators, we can write the killing form in terms of the structure constants of the algebra, with structure constants defined by the commutation relations of the generators. we know tha adjoint representation is given by the structure constants, thus we can write the killing form in terms of the structure constants. this allows us to compute the killing form explicitly for any given lie algebra, once we know its structure constants. Thus it is written in the same basis.

\[
    K_{mn} = \dots
\]

in particular for \(\mathfrak{g} = \mathfrak{su}(2)\) in the previous antihermitian basis we have \(K_{mn} = -\delta_{mn}\) and for the hermitian basis for the complexification \(\mathfrak{sl}(2,\mathbb{C})\)  we have \(K_{mn} = \delta_{mn}\).

There is an ambiguity when it comes to subalgebras: there are a priori two pairings we can define on a subalgebra, the one given by the killing form of the original algebra and the one given by the killing form of the subalgebra (which is a Lie algebra iself). in general they are not the same, thus we need to choose which one to use. whats relevant for us is when they are the same: in any case in which the subalgebra is an invariant subspace of the adjoint then we have the same killing form, since the killing form is defined in terms of the adjoint representation. thus we will always choose a subalgebra that is an invariant subspace of the adjoint, so that we can use the same killing form to define an inner product on it.
thus for a semisimple lie algebra, we have a unique killing form, made by the direct sum of the killing forms of the simple components.
another relevant case is when the subalgebra is real (from the complexification) then the real basis can be used to define the complex basis, therefore structure constants and the pairing matrix in the basis are the same, thus the following corollary holds:

\begin{corollary}
    Let \(\mathfrak{h}\) be a real Lie algebra and \(\mathfrak{g} = \mathfrak{h}_{\mathbb{C}}\) its complexification. Then the killing form of \(\mathfrak{h}\) and the killing form of \(\mathfrak{g}\) coincide on \(\mathfrak{h}\):
    \[
        K^{\mathfrak{h}} = K^{\mathfrak{g}} \big|_\mathfrak{h}.
    \]
\end{corollary}

\begin{proposition}
    The Killing form on \(\mathfrak{g}\) satisfies \(\forall\) ...
    \[
        \dots
    \]
    symmetric pairing - invariant
    their sum is zero
\end{proposition}
\begin{proof}
    [...]
\end{proof}

The notion of “invariance” becomes intuitive when we inspect the interplay between the Killing form on \(\mathfrak{g}\) and the adjoint representation of the Lie group \(G\) with algebra \(\mathfrak{g}\). Namely, from Corollary 2.3.3 we have for all \(X,Y,Z \in \mathfrak{g}\):
\[
    \dots
\]
from which we can see that the infinitesimal transformation \(\mathrm{ad}_{\exp(Z)}\) does not change the Killing form of any two elements of the algebra, thus it is invariant under the adjoint action of the group on the algebra. this is a very important property, since it allows us to use the killing form to define an inner product on the algebra, which is invariant under the adjoint action.

\begin{example}[...]
    [...]
\end{example}